\section{The frontier}

\begin{theorem}[Parent-barrier gap holds]
  \label{thm:collatz-gap-holds}
  \lean{Obsmetrologia.Closure.Collatz.parentBarrierGap_holds}
  \notready
  \uses{thm:collatz-gap, thm:collatz-gap-small-oddcount-slice}
  \emph{(Open frontier --- stated, not proved.)} The gap property itself: nontrivial parent
  barriers never represent $C$. Supported by the finite audit
  \texttt{pnpm run collatz:audit:parent-gap}: through parent-prefix length \(28\)
  it enumerates \(1{,}166{,}059\) parent barriers, checks \(1{,}166{,}057\)
  nontrivial cases, and finds zero residue-gap violations, zero positive
  carry-quotient violations, zero cylinder-realization violations, zero
  bad-cylinder violations, and zero mismatches between the proved-equivalent
  faces. This is evidence, not a proof. The residue refinement is the modular gap
  \(D\cdot r-C>0\) with \(r=C\cdot D^{-1}\bmod M\); Lean constructs the inverse
  from B\'ezout and proves that this residue inequality is equivalent to the
  semigroup gap (\cref{thm:collatz-residue-gap},
  \cref{thm:collatz-residue-gap-reduces}). The two-adic refinement names the
  exact quotient \((D\cdot r-C)/M\), and proves that positivity of that quotient
  is the same open frontier (\cref{thm:collatz-two-adic-carry-quotient-positive}).
  Beyond numerical audit, the gap is now proved unconditionally --- for every
  barrier of the shape, not merely checked ones --- on the whole
  $\mathrm{oddCount}\le3$ slice (\cref{thm:collatz-gap-small-oddcount-slice}): the
  unique two-\(O\) shape \(OOE\) and the two three-\(O\) shapes \(OOEO\), \(OOOE\)
  all satisfy the gap by direct computation, closed off by the trace-exponent
  sandwich that pins their length. The general node remains open at
  $\mathrm{oddCount}\ge4$.
\end{theorem}
\begin{theorem}[Parent-barrier residue gap holds]
  \label{thm:collatz-residue-gap-holds}
  \lean{Obsmetrologia.Closure.Collatz.parentBarrierResidueGap_holds}
  \notready
  \uses{thm:collatz-gap-holds, thm:collatz-residue-gap-reduces}
  \emph{(Same open frontier, residue form.)} The parent-barrier frontier transported
  across the proved equivalence \texttt{ParentBarrierGap} \(\leftrightarrow\)
  \texttt{ParentBarrierResidueGap}. This is the modular inequality statement
  \(C<D\cdot r\) for every nontrivial parent barrier. It introduces no second
  hole: its only unproved input is \cref{thm:collatz-gap-holds}.
\end{theorem}
\begin{theorem}[Parent-barrier residue carry gap holds]
  \label{thm:collatz-residue-carry-gap-holds}
  \lean{Obsmetrologia.Closure.Collatz.parentBarrierResidueCarryGap_holds}
  \notready
  \uses{thm:collatz-residue-gap-holds, thm:collatz-residue-gap-carry}
  \emph{(Same open frontier, carry form.)} The parent-barrier frontier transported
  across the proved carry equivalence. Under the residue congruence, \(C<D\cdot r\)
  is the same as \(M\le D\cdot r-C\), with \(M=2^{|p|+1}\); this is the
  high-\(2\)-adic-digit face of the same red node.
\end{theorem}
\begin{theorem}[Every nonterminal seed has a crossing]
  \label{thm:collatz-crossing-holds}
  \lean{Obsmetrologia.Closure.Collatz.crossingSweep_holds}
  \lean{Obsmetrologia.Closure.Collatz.everyNonterminalContracts_below_wide}
  \notready
  \uses{thm:collatz-crossing-firstpassage, thm:collatz-parent-one-lock-supplies-stream-frontier}
  \emph{(Open frontier --- stated, not proved.)} No nonterminal seed runs an infinite
  coefficient-noncontracting barrier; equivalently every $n>2$ has a finite crossing. This is
  the ``no positive infinite barrier stream'' content. It is no longer a second
  Lean \texttt{sorry}; in the one-lock accounting it is absorbed by the named
  implication \texttt{ParentBarrierOneLockCollapse}. The proved receipt
  \cref{thm:collatz-parent-one-lock-supplies-stream-frontier} says that the
  crossing/no-stream condition follows from the parent-barrier gap once that
  named implication is supplied. The Collatz frontier audit writes this guardrail
  to \texttt{sites/docs/coverage/collatz-frontier-audit.\{json,md\}}: no hard
  trust surfaces, plus a checked one-lock proposition/equivalence layer rather
  than a proved one-lock implication. The finite kernel-checked witness range is
  also substantially widened: \texttt{everyNonterminalContracts\_below\_wide}
  proves the exact open-frontier statement
  ($\exists k,\ \texttt{CoeffContracting}(\mathrm{parityWord}\ n\ k)$, no depth
  bound in the conclusion) for every $2<n<50{,}000$, a hundredfold widening of
  the earlier \texttt{everyNonterminalContracts\_below} range ($n<500$). The
  widening uses a near-linear, tail-recursive \texttt{Bool} verifier
  (\texttt{crossingBool}/\texttt{crossingSweep}) carrying the running trace
  coefficients forward one letter at a time instead of recomputing
  \texttt{parityWord}/\texttt{traceOfWord} from scratch at each candidate depth,
  proved sound once against the exact \texttt{CoeffContracting} semantics and
  then decided on the resulting closed \texttt{Bool} computation alone
  (\texttt{crossingSweep\_holds}, plain kernel \texttt{decide}, no
  \texttt{native\_decide}). It strictly supersedes the earlier audit in range
  without touching it; it does not close the open frontier itself.
\end{theorem}
\begin{theorem}[Crossing is no positive infinite barrier stream]
  \label{thm:collatz-crossing-no-infinite-stream}
  \lean{Obsmetrologia.Closure.Collatz.everyNonterminalHasCrossing_iff_noPositiveInfiniteBarrierStream}
  \leanok
  \uses{thm:collatz-crossing-firstpassage}
  The crossing hypothesis is equivalent, constructively, to the absence of positive infinite
  coefficient-barrier streams: every nonterminal seed's stream has a finite contracting break.
  Lean also records the two pointwise directions: a crossing excludes an infinite barrier stream,
  while failure to reach a first crossing makes every finite prefix noncontracting.
\end{theorem}
\begin{theorem}[Crossing reduces to "eventually contracts"]
  \label{thm:collatz-crossing-reduces}
  \lean{Obsmetrologia.Closure.Collatz.everyNonterminalContracts_iff_crossing}
  \leanok
  \uses{thm:collatz-crossing-no-infinite-stream}
  A \emph{first} coefficient crossing exists as soon as \emph{any} contracting prefix does: the
  least contracting length is automatically a first crossing (\texttt{firstCoeffCrossing\_of\_contracting},
  via a from-scratch well-ordering \texttt{exists\_least} since this core lacks \texttt{Nat.find}).
  Hence the sharp form \texttt{EveryNonterminalContracts} --- every $n>2$ has \emph{some}
  contracting parity prefix ($3^{\mathrm{oddCount}}<2^{\mathrm{len}}$) --- is equivalent to
  \texttt{EveryNonterminalHasCrossing}, dropping the all-prior-noncontracting bookkeeping. The hole
  is thus exactly the bare statement ``every nonterminal seed eventually contracts.''
\end{theorem}
\begin{theorem}[Eventual contraction is no positive infinite barrier stream]
  \label{thm:collatz-eventual-contracts-no-infinite-stream}
  \lean{Obsmetrologia.Closure.Collatz.everyNonterminalContracts_iff_noPositiveInfiniteBarrierStream}
  \leanok
  \uses{thm:collatz-crossing-reduces, thm:collatz-crossing-no-infinite-stream}
  The bare eventual-contraction frontier is equivalent to the absence of positive infinite
  coefficient-barrier streams. This is the dynamical obstruction stripped of first-passage
  bookkeeping: a proof by residue-tree, inverse-thread, or symbolic-dynamical machinery must
  prove this statement, not merely rename the crossing hypothesis.
\end{theorem}
\begin{theorem}[No-stream is parent-barrier cylinder extraction]
  \label{thm:collatz-no-stream-parent-cylinder-extraction}
  \lean{Obsmetrologia.Closure.Collatz.EveryNonterminalHasParentBarrierCylinder}
  \lean{Obsmetrologia.Closure.Collatz.everyNonterminalHasParentBarrierCylinder_iff_crossing}
  \lean{Obsmetrologia.Closure.Collatz.everyNonterminalHasParentBarrierCylinder_iff_noPositiveInfiniteBarrierStream}
  \lean{Obsmetrologia.Closure.Collatz.infiniteBarrierStream_excludes_parentBarrierCylinder}
  \lean{Obsmetrologia.Closure.Collatz.CylinderAvoidingBarrierStream}
  \lean{Obsmetrologia.Closure.Collatz.cylinderAvoidingBarrierStream_iff_positive_infiniteBarrierStream}
  \lean{Obsmetrologia.Closure.Collatz.cylinderAvoidingBarrierStream_iff_infinite_and_no_parentCylinder}
  \lean{Obsmetrologia.Closure.Collatz.cylinderAvoidingBarrierStream_excludes_noPositiveInfiniteBarrierStream}
  \lean{Obsmetrologia.Closure.Collatz.noPositiveInfiniteBarrierStream_excludes_cylinderAvoiding}
  \leanok
  \uses{thm:collatz-crossing-no-infinite-stream, thm:collatz-firstpassage-iff}
  The no-stream frontier is exactly finite parent-barrier cylinder extraction:
  every nonterminal seed has a crossing iff it realizes some first-passage
  cylinder \(pE\) with \(p\) a parent barrier. Lean also proves the important
  negative guardrail: a genuine infinite coefficient-barrier stream cannot
  already realize such a cylinder, because \(pE\) is contracting at its last
  letter. Thus the next red move is not extraction of a bad cylinder from an
  infinite stream; it is proving that every positive nonterminal stream reaches
  a parent-barrier cylinder at all. Lean now names this tightened complement as
  \texttt{CylinderAvoidingBarrierStream}: it is exactly a positive infinite
  barrier stream, because parent-cylinder avoidance is forced by infinitude.
  Such a stream refutes \texttt{NoPositiveInfiniteBarrierStream}; conversely,
  the no-stream frontier excludes all cylinder-avoiding streams.
\end{theorem}
\begin{theorem}[The purl: descent contains the crossing]
  \label{thm:collatz-purl}
  \lean{Obsmetrologia.Closure.Collatz.everyNonterminalContracts_of_allFallBelow}
  \leanok
  \uses{thm:collatz-descent-iff, thm:collatz-crossing-reduces, thm:collatz-no-stream-parent-cylinder-extraction}
  The surplus $S=A\cdot n+C-B\cdot n$ is the single gauge: descent $\iff S<0$. Its two faces are
  the coefficient crossing ($A<B$) and the offset gap ($C<(B-A)\cdot n$). Because $C\ge0$, $S<0$
  forces $A<B$ (\texttt{contracting\_of\_descent}): a descending step \emph{is} a contracting step.
  Hence universal descent below self contains the crossing hypothesis
  (\texttt{everyNonterminalContracts\_of\_allFallBelow}). This explains why crossing appears once
  descent is known, but it does not remove the crossing input from the forward reduction: before a
  seed has a first-passage word, \texttt{ParentBarrierGap} has no parent barrier to act on.
\end{theorem}
\begin{theorem}[Accelerated Collatz]
  \label{thm:collatz-accelerated}
  \notready
  \uses{thm:collatz-conditional, thm:collatz-gap-holds, thm:collatz-crossing-holds}
  \emph{(Open frontier --- derived, rests on the two holes.)} Every positive seed reaches the
  terminal cycle. The derivation from \cref{thm:collatz-conditional} is complete; the node is
  red because it depends on exactly \cref{thm:collatz-gap-holds} and
  \cref{thm:collatz-crossing-holds} and nothing else. The semigroup half can now be replaced by
  the sharper residue inequality via \cref{thm:collatz-residue-conditional}; crossing remains a
  separate red alias in the live Lean graph. More sharply,
  \cref{thm:collatz-frontier-inputs-gap-and-one-lock} proves that the whole frontier package is
  the gap together with the named one-lock implication; after assuming the residue gap,
  accelerated termination is equivalent to absence of a positive infinite barrier stream
  (\cref{thm:collatz-residue-exact-obstruction}).
\end{theorem}
\begin{theorem}[Residue-frontier accelerated Collatz]
  \label{thm:collatz-residue-frontier-accelerated}
  \notready
  \uses{thm:collatz-residue-gap-holds, thm:collatz-crossing-holds, thm:collatz-residue-conditional}
  The same headline routed through the modular frontier node:
  \texttt{parentBarrierResidueGap\_holds} plus the crossing frontier implies
  accelerated Collatz. This theorem has no new mathematical content beyond
  \cref{thm:collatz-accelerated}; it makes the visible arithmetic red node the
  modular inequality \(C<D\cdot r\).
\end{theorem}

\begin{theorem}[The arithmetic zeta-star finite bridge]
  \label{thm:arithmetic-zeta-star-finite-bridge}
  \lean{Obsmetrologia.AperiodicOrder.Tiling.zetaStar_bridge_evidence}
  \leanok
  \uses{thm:collatz-affine-trace, thm:finite-displacement-delone, thm:finite-diffraction-mass, thm:salem-pisot-regime-split}
  The finite zeta-star object is the affine-trace residual read on the finite
  displacement carrier \(\mathrm{Fin}(N+1)\).  Lean proves that the residual
  vanishes at every displacement, so its zero-set is the full finite carrier;
  that zero-set is Delone in the finite displacement sense of
  \texttt{Diffraction.lean}, has total finite diffraction mass \(N+1\), and its
  finite pure-point readout is exactly the Pisot side of the existing StepN
  phase classifier.  The bundled bridge also exposes the Pisot-side finite
  diffraction readout as a positive pure-point atom and a finite pure-point
  quasicrystal surrogate.  This is a finite carrier theorem, not an analytic
  zeta function.
\end{theorem}

\begin{theorem}[The zeta-star arithmetic settling certificate]
  \label{thm:zeta-star-arithmetic-settling-certificate}
  \lean{Obsmetrologia.AperiodicOrder.Tiling.zetaStar_arithmeticSettling_certificate}
  \leanok
  \uses{thm:arithmetic-zeta-star-finite-bridge, thm:collatz-saturating-depth, thm:stepn-chu-cayley-hamilton, thm:finite-predictor-phase-type}
  Lean bundles the finite zeta-star bridge as an arithmetic-settling
  certificate: affine-trace zeros fill the finite displacement carrier, the
  affine-trace identity holds at every finite displacement with denominator
  coefficient \(2^d\), the zero-set is finite-discrete Delone, the total finite
  diffraction mass is \(N+1\), odd seeds carry the 2-adic shortcut valuation
  witness, the retained Fibonacci substitution matrix is primitive and
  satisfies its characteristic polynomial, and the finite Fourier-side readout
  agrees with
  \(\texttt{phaseType}\).  This is the strongest statement currently proved that
  ζ-star is over the corpus's Collatz/2-adic/substitution machinery.
\end{theorem}

\begin{theorem}[The zeta-star finite mass has Pisot pure-point readout]
  \label{thm:zeta-star-finite-mass-pure-point}
  \lean{Obsmetrologia.AperiodicOrder.Tiling.zetaStar_finiteDiffractionMass_purePoint_of_pisot}
  \leanok
  \uses{thm:zeta-star-arithmetic-settling-certificate, thm:salem-pisot-regime-split}
  On the Pisot side of the existing StepN phase classifier, Lean packages
  \(\mathrm{finiteDiffractionMass}(Z(\zeta^\star))\) as the atom mass of a
  finite Fourier-side readout whose component is pure point.  The theorem also
  records that this atom mass is positive and exactly the finite diffraction
  mass of the zeta-star zero-set.  This is the pure-point
  statement currently proved by the finite carrier machinery.
\end{theorem}

\begin{theorem}[There exists a zeta-star with pure-point finite mass]
  \label{thm:exists-zeta-star-finite-mass-pure-point}
  \lean{Obsmetrologia.AperiodicOrder.Tiling.exists_zetaStar_finiteDiffractionMass_purePoint}
  \leanok
  \uses{thm:zeta-star-finite-mass-pure-point, thm:zeta-star-arithmetic-settling-certificate}
  For every finite displacement window \(N\), Lean constructs the canonical
  zeta-star witness and proves the finite form of the pure-point diffraction
  claim: its finite phase is the consistent face, its finite Fourier-side
  component is pure-point, and the readout atom is exactly
  \(\mathrm{finiteDiffractionMass}(Z(\zeta^\star))\) with positive mass.
\end{theorem}

\begin{theorem}[The zeta-star zero-set is a finite pure-point quasicrystal on the Pisot side]
  \label{thm:zeta-star-finite-pure-point-quasicrystal}
  \lean{Obsmetrologia.AperiodicOrder.Tiling.zetaStar_finitePurePointQuasicrystal_of_pisot}
  \leanok
  \uses{thm:exists-zeta-star-finite-mass-pure-point, thm:zeta-star-arithmetic-settling-certificate, thm:finite-pure-point-quasicrystal}
  On the Pisot side, \(Z(\zeta^\star)\) instantiates the finite quasicrystal
  carrier now defined in \texttt{Diffraction.lean}: its zero-set is
  finite-displacement Delone, its finite Fourier-side readout is pure-point,
  and its atom mass is the total finite diffraction mass.  This is the
  finite surrogate proved for the requested quasicrystal statement.
\end{theorem}

\begin{theorem}[There exists a finite zeta-star quasicrystal witness]
  \label{thm:exists-zeta-star-finite-quasicrystal}
  \lean{Obsmetrologia.AperiodicOrder.Tiling.exists_zetaStar_finitePurePointQuasicrystal}
  \leanok
  \uses{thm:zeta-star-arithmetic-settling-certificate, thm:zeta-star-finite-pure-point-quasicrystal}
  For every finite displacement window \(N\), Lean constructs a concrete
  zeta-star object: the canonical witness has Collatz seed \(1\) and the Pisot
  StepN side.  It carries the arithmetic-settling certificate and instantiates
  the finite pure-point quasicrystal surrogate over \(Z(\zeta^\star)\).  This is
  the finite existential theorem the corpus currently proves.
\end{theorem}

\begin{theorem}[The bundled finite Dyson quasicrystal form of zeta-star]
  \label{thm:zeta-star-finite-dyson-quasicrystal-form}
  \lean{Obsmetrologia.AperiodicOrder.Tiling.exists_zetaStar_finiteDysonQuasicrystalForm}
  \leanok
  \uses{thm:exists-zeta-star-finite-mass-pure-point, thm:exists-zeta-star-finite-quasicrystal, thm:zeta-star-arithmetic-settling-certificate}
  Lean packages the single finite diffraction statement currently proved:
  there exists a zeta-star witness whose arithmetic certificate is present,
  whose zero-set is Delone on the finite displacement carrier, whose finite
  diffraction mass is \(N+1\), whose phase lies in the bigraded trichotomy and
  is the consistent face, whose finite Fourier-side readout is pure-point with
  positive atom mass exactly equal to \(\mathrm{finiteDiffractionMass}
  (Z(\zeta^\star))\), and whose zero-set inhabits the finite pure-point
  quasicrystal carrier.
\end{theorem}

\begin{theorem}[The zeta-star coalgebraic promotion apparatus]
  \label{thm:zeta-star-coalgebraic-promotion-apparatus}
  \lean{Obsmetrologia.AperiodicOrder.Tiling.zetaStar_coalgebraicPromotion_apparatus}
  \leanok
  \uses{thm:zeta-star-finite-dyson-quasicrystal-form, thm:metallic-power-dichotomy, thm:coalgebra-dual-of-algebra, thm:finite-scott-continuous, thm:coalgebraic-crossbase-mld}
  Lean packages the finite zeta-star form together with the existing
  coalgebraic machinery: the metallization dichotomy supplies the
  finite-algebra/limit-coalgebra boundary, finite directed joins are attained,
  finite charge closures are Scott-continuous, settling already has the
  algebra/coalgebra duality package, and up-to-closure coinduction classifies
  geometric MLD by settled equality.  This is the current internal apparatus
  for the analytic promotion; it does not yet instantiate the zeta-star
  displacement coalgebra.
\end{theorem}

\begin{theorem}[The zeta-star algebra/coalgebra boundary]
  \label{thm:zeta-star-algebra-coalgebra-boundary}
  \lean{Obsmetrologia.AperiodicOrder.Tiling.zetaStar_algebraCoalgebraBoundary}
  \leanok
  \uses{thm:zeta-star-coalgebraic-promotion-apparatus, thm:metallic-power-dichotomy}
  Lean records the reframe explicitly: the finite displacement windows carry
  the zeta-star finite Dyson form, the already-proved metallization dichotomy is
  the finite algebra versus limit coalgebra boundary, and the Scott/coinductive
  apparatus is the analytic side.  This proves the ``finite = algebra,
  analytic = coalgebra'' correspondence so far; the remaining red work is to build the actual
  zeta-star displacement coalgebra and prove its diffraction theorem.
\end{theorem}

\begin{theorem}[The finite observed zeta-star displacement coalgebra]
  \label{thm:zeta-star-finite-observed-displacement-coalgebra}
  \lean{Obsmetrologia.AperiodicOrder.Tiling.canonicalZetaStarDisplacementCoalgebra_quasicrystal}
  \leanok
  \uses{thm:zeta-star-finite-dyson-quasicrystal-form, thm:zeta-star-algebra-coalgebra-boundary}
  Lean constructs the finite observed displacement coalgebra proved so far:
  its state carrier is the finite zeta-star state space, its observation is the
  tuple consisting of zero-set, finite diffraction mass, finite component, and
  phase readout, and its finite unfold step preserves that observation.  At the
  canonical seed, this coalgebra carries the finite pure-point quasicrystal
  witness.  This is still the idempotent finite carrier, not the analytic
  displacement coalgebra.
\end{theorem}

\begin{theorem}[The zeta-star analytic-richness recovery interface]
  \label{thm:zeta-star-analytic-richness-recovery}
  \lean{Obsmetrologia.AperiodicOrder.Tiling.zetaStar_analyticRichnessRecovery}
  \leanok
  \uses{thm:zeta-star-finite-observed-displacement-coalgebra, thm:observed-upto-closure-coinduction}
  Lean identifies why the coalgebraic bridge was losing analytic content: the
  bare settled-value kernel classified MLD but did not carry an observation.
  The recovered interface uses observed up-to-closure coinduction, so any
  diffraction, phase, spectral, or mass readout that is proved compatible with
  closure is preserved by the coalgebra.  It now includes the finite observed
  zeta-star displacement coalgebra as the carrier constructed so far.  The remaining
  zeta-star gap is the non-idempotent analytic displacement coalgebra and the
  observed-closure compatibility of its diffraction observation, not an unnamed
  analytic fog.
\end{theorem}

\begin{theorem}[There exists a finite zeta-star arithmetic witness with odd shortcut]
  \label{thm:exists-zeta-star-arithmetic-odd-shortcut}
  \lean{Obsmetrologia.AperiodicOrder.Tiling.exists_zetaStar_arithmeticSettling_with_odd_shortcut}
  \leanok
  \uses{thm:zeta-star-analytic-richness-recovery, thm:zeta-star-arithmetic-settling-certificate}
  The canonical finite witness can be read with all currently proved arithmetic
  facts visible at once: seed \(1\), Pisot StepN phase, the arithmetic
  settling certificate, the finite pure-point quasicrystal surrogate, and the
  explicit odd-seed 2-adic shortcut valuation for \(3\cdot\texttt{seed}+1\).
  This sharpens the finite existential bridge; it is still not the analytic
  Delone or Dyson/RH theorem.
\end{theorem}

\begin{theorem}[The zeta-star finite readout agrees with phaseType]
  \label{thm:zeta-star-finite-phase-type}
  \lean{Obsmetrologia.AperiodicOrder.Tiling.zetaStar_phaseType_readout}
  \leanok
  \uses{thm:exists-zeta-star-arithmetic-odd-shortcut, thm:zeta-star-finite-pure-point-quasicrystal, thm:zeta-star-arithmetic-settling-certificate, thm:zeta-star-finite-mass-pure-point, thm:finite-predictor-phase-type}
  The finite Fourier-side readout of \(Z(\zeta^\star)\) lands in the corpus's
  existing finite predictor-level \(\texttt{phaseType}\) map: applying
  \(\texttt{lebesgueType}\) to the zeta-star readout component gives exactly
  \(\texttt{zetaStarPhaseType}\).  Lean also proves the companion equivalence
  \(\texttt{zetaStar\_readout\_purePoint\_iff\_phaseType\_consistent}\):
  finite pure-point readout is the same as the consistent/Pisot phase.  The
  certificate also records the zeta-star phase's entry into the corpus's
  bigraded trichotomy.
\end{theorem}

\begin{theorem}[The zeta-star zero-set is analytically Delone]
  \label{thm:zeta-star-analytic-delone}
  \notready
  \uses{thm:zeta-star-coalgebraic-promotion-apparatus, thm:arithmetic-zeta-star-finite-bridge, thm:zeta-star-finite-pure-point-quasicrystal, thm:zeta-star-finite-phase-type}
  \emph{(Open coalgebraic/Scott frontier --- stated, not formalized.)} Promote
  the finite zero-set approximants of \(\zeta^\star\) from the algebraic carrier
  \(\mathrm{Fin}(N+1)\) to the directed-join displacement coalgebra, then prove
  the Delone conditions there.  This is the dual fixpoint problem: the finite
  side is the initial-algebra approximation, while the analytic side is the
  final-coalgebra/Scott limit.  The current Lean theorem supplies only the
  finite-discrete Delone surrogate and its bundled finite diffraction form.
\end{theorem}

\begin{theorem}[Dyson/RH as pure-point diffraction of \(Z(\zeta^\star)\)]
  \label{thm:zeta-star-rh-pure-point}
  \notready
  \uses{thm:zeta-star-analytic-delone, thm:zeta-star-finite-phase-type, thm:salem-color-line-fourier-atom}
  \emph{(Open coalgebraic/Scott frontier --- stated, not formalized.)} The
  intended RH form is that the diffraction measure of the coalgebraic
  directed-join zero-set \(Z(\zeta^\star)\) is pure point in the corpus's
  quasicrystal sense.  The finite bridge shows where the statement attaches to
  Collatz affine traces, finite diffraction mass, and the Pisot/Salem phase
  readout; the missing step is the Scott-continuous promotion from finite
  approximants to the final-coalgebra displacement object.
\end{theorem}
